\section{Context construction is runtime scheduling}\label{sec:C20:context}
Context management determines which parts of execution state the policy can use. A journal may be complete while $C_t$ omits the operation that explains a pending effect. Conversely, a summary may assert completion that the journal never committed. The harness therefore needs a projection contract between authoritative state and model-visible state. For a projection claiming effect-aware planning, the relevant unresolved operations and active constraints need a usable representation. Exact authority and sensitive effect evidence can remain privileged, with a policy-safe status projection; copying every grant or credential into the prompt is unnecessary.

A projection can be tested against invariants without requiring an exact semantic equivalence oracle. For example, inject a pending-operation record and verify that the constructed context includes its stable identifier and uncertainty state. Insert a revoked capability and check that the tool registry removes it regardless of whether the model remembers an earlier grant. These checks establish structural properties. Whether a compressed explanation preserves every decision-relevant fact remains a behavioral question and requires task-level evaluation.

Agentic Context Management supplies model-directed offloading and retrieval mechanisms~\cite{corpus260723809}. Its live-run raw mirror can preserve access to discarded content, but archival retention does not guarantee that the agent retrieves the right item at the right time. Its reported benefits also depend on a capable backbone and reimplemented baselines. This distinction matters for runtime design: storing raw history is a recovery affordance; selecting useful context is an additional policy with cost, latency, and failure modes. The snapshot-only restore path discussed in Chapter~\ref{ch:C13} does not preserve a previously separate raw archive by itself~\cite{repo_acm2026}.

An execution scheduler can reserve a protected context region for authoritative status and separately allocate a discretionary region for retrieved evidence. Protection should mean application-controlled serialization, not merely an instruction asking the model to respect a textual heading. The trusted region should be generated from current state on every call. If a retrieved document can overwrite it through string concatenation or malformed role boundaries, the separation is nominal rather than enforced.

The resource constraint couples context and scheduling. Long tool outputs increase prefill cost and may delay other sessions. An aggressive truncation policy improves throughput but can remove the evidence needed for reconciliation. A token budget should therefore distinguish required status tokens, candidate evidence tokens, and speculative reasoning allocations. When the budget cannot accommodate required state, the runtime should reduce work, retrieve narrower evidence, or request a bounded continuation. Silent removal of constraints is not a legitimate optimization.

Evaluate these choices with task-length and output-size strata. Average success can hide a regime where compaction helps short tasks but fails after the first interruption. Include a baseline that keeps authoritative status fixed while varying only discretionary history. This separates failures of the projection contract from failures of the model to interpret otherwise available evidence. Report the cost of reconstructing omitted evidence, not only the reduction in visible prompt length.

